\chapter{1997年上海卷（文）}
\section{单选题}

\begin{problem}
设全集是实数集 \(\R\)，\(M = \{x \mid x \leqslant 1 + \sqrt{2}, x \in \R\}\)，\(N = \{1, 2, 3, 4\}\)，则 \(\overline{M} \cap N\) 等于（\quad）

\choices
    {\(\left|4\right| \)}
    {\(\{3,4\} \)}
    {\(\{2,3,4\} \)}
    {\(\{1,2,3,4\} \)}

\end{problem}

\begin{problem}
三个数 \(6^{0.7}, 0.7^{6}, \log_{0.76}\) 的大小顺序是（\quad）

\choices
    {\(0.7^{6} < \log_{0.76} < 6^{0.7}\)}
    {\(0.7^{6} < 6^{0.7} < \log_{0.76}\)}
    {\(\log_{0.76}<6^{0.7}<0.7^{6} \)}
    {\(\log_{0.76}<0.7^{6}<6^{0.7} \)}

\end{problem}

\begin{problem}
在下列命题中，假命题是（\quad）

\choices
    {若 \(a\)、\(b\) 是异面直线，则一定存在平面 \(\alpha\) 过 \(a\) 且与 \(b\) 平行}
    {若 \(a\)、\(b\) 是异面直线，则一定存在平面 \(\alpha\) 过 \(a\) 且与 \(b\) 垂直}
    {若 \(a\)、\(b\) 是异面直线，则一定存在平面 \(\alpha\) 与 \(a\)、\(b\) 所成角相等}
    {若 \(a\)、\(b\) 是异面直线，则一定存在平面 \(\alpha\) 与 \(a\)、\(b\) 的距离相等}

\end{problem}

\begin{problem}
设 \(k > 1\)，则关于 \(x, y\) 的方程 \((1 - k)x^{2} + y^{2} = k^{2} - 1\) 所表示的曲线是（\quad）

\choices
    {长轴在 \(y\) 轴上的椭圆}
    {长轴在 \(x\) 轴上的椭圆}
    {实轴在 \(y\) 轴上的双曲线}
    {实轴在 \(x\) 轴上的双曲线}

\end{problem}

\begin{problem}
如果直线 \(l\) 沿 \(x\) 轴负方向平移 3 个单位，再沿 \(y\) 轴正方向平移 1 个单位后，又回到原来的位置，那么直线 \(l\) 的斜率是（\quad）

\choices
    {\(-\frac{1}{3}\)}
    {\(-3\)}
    {\(\frac{1}{3} \)}
    {\(3\)}

\end{problem}

\begin{problem}
设 \(f(n) = 1 + \frac{1}{2} + \frac{1}{3} + \cdots + \frac{1}{3n-1} (n \in N) \)，则 \(f(n+1) - f(n) \) 等于（\quad）

\choices
    {\(\frac{1}{3n+2} \)}
    {\(\frac{1}{3n} + \frac{1}{3n + 1}\)}
    {\(\frac{1}{3n + 1} +\frac{1}{3n + 2}\)}
    {\(\frac{1}{3n} + \frac{1}{3n + 1} + \frac{1}{3n + 2}\)}

\end{problem}

\section{填空题}

\begin{problem}
方程 \(\lg(1-3x)=\lg(3-x)+\lg(7+x) \) 的解是 \fillinblank{}.

\end{problem}

\begin{problem}
方程 \(\sin2x=\frac{1}{2} \) 在 \([-2\pi,2\pi] \) 内解的个数是 \fillinblank{}.

\end{problem}

\begin{problem}
已知 \(a = \frac{-3 - \i}{1 + 2\i} \)（\(\i\) 是虚数单位），那么 \(a^{2} = \) \fillinblank{}.

\end{problem}

\begin{problem}
设圆 \(x^{2} + y^{2} - 4x - 5 = 0 \) 的弦 \(AB\) 的中点为 \(P(3,1) \) ，则直线 \(AB\) 的方程是 \fillinblank{}.

\end{problem}

\begin{problem}
若 \((3x+1)^{n}(n\in N) \) 的展开式中各项系数的和是256，则展开式中 \(x^{2} \) 的系数是 \fillinblank{}.

\end{problem}

\begin{problem}
函数 \(f(x)=3\sin x\cos x-1 \) 的最大值是 \fillinblank{}.

\end{problem}

\section{选作题}

\begin{problem}
设 \(0 < a < b\)，则 \(\displaystyle\lim _{n \rightarrow \infty} \frac {4 b ^{n}}{a ^{n} - b ^{n}}=\)\fillinblank{}.

\end{problem}

\begin{problem}
设正四棱锥底面边长为 \(4 \, cm \) ，侧面和底面所成的二面角是 \(60^{\circ} \) ，则这个棱锥的侧面积是 \(cm^{2} \) \fillinblank{}.

\end{problem}

\begin{problem}
圆柱形容器的内壁底半径为 5cm，两个直径为 5cm 的玻璃小球都浸没于容器的水中. 若取出这两个小球，则容器内的水面将下降 cm \fillinblank{}.

\end{problem}

\begin{problem}
从集合 \(\{0,1,2,3,5,7,11\}\) 中任取3个元素分别作为直线方程 \(Ax + By + C = 0\) 中的 \(A, B, C\)，所得经过坐标原点的直线有 条. （结果用数值表示）\fillinblank{}

\end{problem}

\begin{problem}
\(\displaystyle\lim_{n\to\infty}(1-\frac{2}{n})-2n= \) \fillinblank{}.

\end{problem}

\begin{problem}
设 \(\bs{a}=(m+1)\bs{i}-3\bs{j},\bs{b}=\bs{i}+(m-1)\bs{j},(\bs{a}+\bs{b})\perp(\bs{a}-\bs{b}) \) ，则 \(m =\) \fillinblank{}

\end{problem}

\begin{problem}
设 \(f(x)=x^{2}(2-x) \) ，则 \(f(x) \) 的单调递增区间是 \fillinblank{}.

\end{problem}

\begin{problem}
从集合 \(\{0,1,2,3,5,7,11\}\) 中任取3个元素分别作为直线方程 \(Ax + By + C = 0\) 中的 \(A,B,C\) ，所得恰好经过坐标原点的直线的概率是 \fillinblank{}

\end{problem}

\section{解答题}

\begin{problem}
已知 \(\mathrm{tg}\frac{\alpha}{2} = \frac{1}{2}\) 求 \(\sin \left(\alpha +\frac{\pi}{6}\right)\) 的值. 解：

\end{problem}

\begin{problem}
公园要建造一个圆形的喷水池. 在水池中央垂直于水面安装一个花形柱子 \(OA\)，\(O\) 恰在水面中心，\(OA = 1.25\) 米，安置在柱子顶端 \(A\) 处的喷头向外喷水，水流在各个方向上沿形状相同的抛物线路径落

下，且在过 \(OA\) 的任一平面上抛物线路径如图所示. 为使水流形状较为漂亮，设计成水流在到 \(OA\) 距离为 1 米处达到距水面最大高度 2.25 米. 如果不计其他因素，那么水池的半径至少要多少米，才能使喷出的水流不致落到池外\(\perp\) 解：

\end{problem}

\begin{problem}
如图，在三棱柱 \(ABC - A'B'C'\) 中，四边形 \(A'ABB'\) 是菱形，四边形 \(BCC'B'\) 是矩形，\(C'B' \perp AB\). （1）求证: 平面 \(CA'B \perp A'AB\)；（2）若 \(C'B' = 3\)，\(AB = 4\)，\(\angle ABB' = 60^{\circ}\)，

求 \(AC'\) 与平面 \(BCC'\) 所成角的大小. （用反三角函数表示） [证] （1）[解] （2）

\end{problem}

\begin{problem}
设虚数 \(z_{1}\)、\(z_{2}\) 满足 \(z_{1}^{2}=z_{2}\). （1）若 \(z_{1}\)、\(z_{2}\) 又是一个实系数一元二次方程的两个根，求 \(z_{1}\)、\(z_{2}\)；（2）若 \(z_{1}=1+m\i\)（\(\i\) 为虚数单位），\(|z_{1}|\leqslant\sqrt{2}\)，复数 \(\omega=z_{2}+3\)，求 \(|\omega|\) 的取值范围.

\end{problem}

\begin{problem}
如图，抛物线 \(C_{1}: y = -x^{2} \) 与抛物线 \(C_{2}: y = x^{2} - 2ax (a > 0) \) 交于 \(O\)、\(A\) 两点.

（1）把 \(C_{1}\) 与 \(C_{2}\) 所围成的图形（阴影部分）绕 \(x\) 轴旋转一周，求所得几何体的体积 \(V\)；（2）若过原点的直线 \(l\) 与抛物线 \(C_{2}\) 所围成的图形面积为 \(\frac{9}{2}a^{3}\)，求直线 \(l\) 的方程. [解] （1）[解] （2）

\end{problem}

\begin{problem}
如图，抛物线方程为 \(y^{2}=p(x+1)(p>0)\)，直线 \(x+y=m\) 与 \(x\) 轴的交点在抛物线的准线的右边. （1）求证：直线与抛物线总有两个交点；（2）设直线与抛物线的交点为 \(Q, R\)，\(OQ \perp OR\)，求 \(p\) 关于 \(m\) 的函数 \(f(m)\) 的表达式；（3）在（2）的条件下，若抛物线焦点 \(F\) 到直线 \(x + y = m\) 的距离为 \(\frac{\sqrt{2}}{2}\)，求此直线的方程. [证] （1）

[解] （2）[解] （3）

\end{problem}

\begin{problem}
设数列 \(\{a_{n}\}\) 的首项 \(a_{1}=1\)，前 \(n\) 项和 \(S_{n}\) 满足关系式：\(3 t S _{n} - (2 t + 3) S _{n - 1} = 3 t \quad (t > 0, n = 2, 3, 4, \dots)\). （1）求证：数列 \(\{a_{n}\}\) 是等比数列；（2）设数列 \(\{a_{n}\}\) 的公比为 \(f(t)\)，作数列 \(\{b_{n}\}\)，使 \(b_{1}=1\)，\(b_{n}=f\left(\frac{1}{b_{n-1}}\right)(n=2,3,4,\cdots)\)，求 \(b_{n}\)；（3）求和：\(b_{1}b_{2}-b_{2}b_{3}+b_{3}b_{4}-\cdots+b_{2n-1}b_{2n}-b_{2n}b_{2n+1}\) [证] （1）[解] （2）[解] （3）
\end{problem}

